\section{Results and discussion}\label{sec:results}

Wall wettability affects both the nucleation onset and the subsequent bubble growth \cite{kim2015,chen2019}. This section examines the effect of wettability on the nucleation process using our molecular dynamics simulations.

Fig.~\ref{fig:onset} shows the temporal evolution of the vapor fraction near the wall for contact angles of 35$^\circ$, 70$^\circ$ and 110$^\circ$. At a wall temperature of 300 K, the onset time of nucleation decreased from 350 ps on the 110$^\circ$ surface to 125 ps on the 35$^\circ$ surface, which indicates that a smaller contact angle accelerates nucleation. The 70$^\circ$ surface exhibited an intermediate onset time of 210 ps. These results are consistent with a monotonic variation of the onset time with the contact angle over the range studied.

\begin{equation}
  t_{\mathrm{on}} = t_0 \exp\left(\frac{\Delta G^*}{k_B T_w}\right)
  \label{eq:onset}
\end{equation}

The increase in nucleation onset time with contact angle can be interpreted using classical nucleation theory, as expressed in Eq.~\eqref{eq:onset}. Within this framework, the trend may be attributed to a lower energy barrier $\Delta G^*$ on more hydrophilic surfaces, resulting from the enhanced solid--liquid interaction \cite{carey2018}. This result suggests that interfacial wettability affects the onset of phase change and that, under the conditions studied here, wettability could be considered in the design of cooling surfaces for electronic devices.
